\section{Positional Geometry and Learned Use}
\label{sec:compact-geometry}\label{sec:geometry-proofs}
\subsection{Cross-Gram formula, invariance and effective rank}
\paragraph{Closed-form cross-Gram.} For $\Delta\sim\mathrm{Unif}[0,L]$ and
$x_\omega(\Delta)=[\cos(\omega\Delta)\ \ \sin(\omega\Delta)]$, write
$d=(\omega{-}\nu)L$, $s=(\omega{+}\nu)L$, $a(t)=\sin t/t$ and
$b_\star(t)=(1{-}\cos t)/t$. Expanding the products of sinusoids and integrating
termwise gives
\begin{equation}
H_{\omega\nu}=\mathbb E[x_\omega^\top x_\nu]=\frac12
\begin{bmatrix}
a(d)+a(s) & b_\star(s)-b_\star(d)\\
b_\star(s)+b_\star(d) & a(d)-a(s)
\end{bmatrix},
\label{eq:cross-gram}
\end{equation}
and $S_\omega=H_{\omega\omega}$. Since $Q_{\omega\nu}=S_\omega^{-1/2}H_{\omega\nu}S_\nu^{-1/2}$
is the whitened cross-Gram of two two-dimensional subspaces, its singular values
are the canonical correlations, and $c_{\omega\nu}=\frac12\|Q_{\omega\nu}\|_F^2$
is invariant under $x_\omega\mapsto x_\omega O$ for any invertible $O$ (in
particular under rotation of the phase inside a pair), because $S_\omega$
transforms congruently and the whitening cancels $O$.

\paragraph{Basis invariance.}
The cancellation can be stated exactly. Under independent invertible basis
changes $x_\omega\mapsto x_\omega A_\omega$ and
$x_\nu\mapsto x_\nu A_\nu$, let
$U_j=S_j^{1/2}A_j(A_j^\top S_j A_j)^{-1/2}$ for $j\in\{\omega,\nu\}$. Then
$U_j^\top U_j=I$, and the transformed whitened cross-Gram is
$\widetilde Q_{\omega\nu}=U_\omega^\top Q_{\omega\nu}U_\nu$. Left and right
orthogonal factors preserve the two canonical correlations and therefore
$c_{\omega\nu}$. Thus $c_{\omega\nu}$ is an intrinsic measure of overlap
between the two positional subspaces.

\begin{theorem}[Spectral budget identity]
\label{thm:budget}
For the block-whitened Gram matrix $\Gamma$ with diagonal blocks $I_2$,
$\operatorname{tr}\Gamma=2K$ and
$\operatorname{tr}(\Gamma^2)=2K[1+(K-1)\bar c]$. Hence its
R\'enyi-$2$ effective rank is given by \eqref{eq:budget-identity}.
\end{theorem}

\begin{proof}[Proof of Theorem~\ref{thm:budget}]
$\Gamma$ is the block matrix with blocks $\Gamma_{ij}=Q_{\omega_i\omega_j}$ and
$\Gamma_{ii}=I_2$. Hence $\operatorname{tr}\Gamma=\sum_{i=1}^K\operatorname{tr}I_2=2K$.
For the second moment,
$\operatorname{tr}(\Gamma^2)=\sum_{i,j}\operatorname{tr}(\Gamma_{ij}\Gamma_{ji})
=\sum_i\operatorname{tr}(I_2)+\sum_{i\neq j}\|Q_{\omega_i\omega_j}\|_F^2
=2K+2K(K-1)\bar c$, using $\Gamma_{ji}=\Gamma_{ij}^\top$ and the definition
$\bar c=\frac{1}{K(K-1)}\sum_{i\neq j}c_{\omega_i\omega_j}$ with
$\|Q_{ij}\|_F^2=2c_{ij}$. Dividing gives \eqref{eq:budget-identity}.
\end{proof}

The identity determines R\'enyi-$2$ effective rank exactly from the mean pairwise
redundancy. Shannon effective rank and log-determinant retain higher-order
multi-subspace dependence on the full spectrum.

\begin{proof}[Proof of Proposition~\ref{prop:collapse}]
Put $t=\Delta/L$, $x=\omega L$, and use $\theta=x^2$ as the local parameter.
The rescaled pair has the analytic basis
\begin{align}
u_1(\theta,t)&=\cos(\sqrt\theta\,t)
 =1-\tfrac12\theta t^2+O(\theta^2),\\
u_2(\theta,t)&=\frac{\sin(\sqrt\theta\,t)}{\sqrt\theta}
 =t-\tfrac16\theta t^3+O(\theta^2).
\label{eq:slow-basis-expansion}
\end{align}
Thus $V_\omega\to V_0=\operatorname{span}\{1,\Delta\}$.  The fourth-order
coefficient follows by retaining the part of the first perturbation transverse
to $V_0$. Under the
uniform measure on $t\in[0,1]$,
\begin{align}
\Pi_0^\perp t^2&=p_2(t):=t^2-t+\tfrac16,\\
\Pi_0^\perp t^3&=p_3(t):=t^3-\tfrac9{10}t+\tfrac15,
\label{eq:slow-residual-polynomials}
\end{align}
where $\Pi_0^\perp$ is the $L_2[0,1]$ projection off
$\operatorname{span}\{1,t\}$.  In the base $[1,t]$, its Gram and the residual
perturbation Gram are
\begin{equation}
G_0=
\begin{bmatrix}1&1/2\\[1pt]1/2&1/3\end{bmatrix},\qquad
M=
\begin{bmatrix}
\langle-p_2/2,-p_2/2\rangle&\langle-p_2/2,-p_3/6\rangle\\
\langle-p_3/6,-p_2/2\rangle&\langle-p_3/6,-p_3/6\rangle
\end{bmatrix}
=\begin{bmatrix}1/720&1/1440\\[1pt]1/1440&1/2800\end{bmatrix}.
\label{eq:slow-residual-gram}
\end{equation}
The squared chordal distance between the two nearby planes is therefore
\begin{equation}
2-\lVert Q_{x,y}\rVert_F^2
=(x^2-y^2)^2\operatorname{tr}(G_0^{-1}M)+O(\epsilon^6)
=\frac{19}{12600}(x^2-y^2)^2+O(\epsilon^6),
\label{eq:slow-collapse-expanded}
\end{equation}
with $\epsilon=\max\{|x|,|y|\}\to0$.  The exact coefficient is reproduced
with rational polynomial inner products by
\texttt{scripts/analysis/full\_rope\_collision\_audit.py}; at
$x{=}0.05,y{=}0.10$ its analytic Gram gives exact/leading ratio $1.00058$.

For the softmax metric with fixed $p$,
$F=\operatorname{diag}(p)-pp^\top$ satisfies $F\mathbf1=0$, so the constant
direction is annihilated before whitening.  Writing
$\overline h=h-\mathbb E_p h$, the same expansion gives
\begin{equation}
\frac{\overline{\sin(\omega\Delta)}}{\omega}
 \longrightarrow \Delta-\mathbb E_p\Delta,
\qquad
-\frac{2\overline{\cos(\omega\Delta)}}{\omega^2}
 \longrightarrow \Delta^2-\mathbb E_p\Delta^2.
\end{equation}
These functions span the centred limit whenever $p$ has nondegenerate support
on at least three distances. The separation weights $p$ therefore determine
which polynomial directions survive centering.
\end{proof}

\subsection{Finite-window overlap in the identification experiment}
\label{sec:finite-window-rank}

The slow-limit proposition has a finite-window counterpart that can be calculated directly on the identification grid. Set $b=256$, $K=32$, $L=256$, and retain the $n$ slowest frequencies of $\omega_j=b^{-j/K}$. Their dimensionless frequencies are $\omega_jL=1.1892,1.4142,1.6818,2,\ldots$ in slow-to-fast order. The smallest is already greater than one, so we evaluate the full cross-Gram in~\eqref{eq:cross-gram} without a small-phase approximation.

\begin{table}[ht]
\centering\small
\caption{\textbf{Finite-window positional concentration on the paired-training grid.} The separation measure is uniform on $[0,L]$. The nominal coordinate count is $2n$; $r_2$ is the effective rank of the block-whitened Gram.}
\label{tab:finite-window-rank}
\begin{tabular}{@{}rrrr@{}}
\toprule
Slowest pairs $n$ & Largest $\omega L$ & Coordinates & $r_2$ \\
\midrule
2 & 1.41 & 4 & 2.00030 \\
4 & 2.00 & 8 & 2.00361 \\
6 & 2.83 & 12 & 2.02263 \\
8 & 4.00 & 16 & 2.11474 \\
12 & 8.00 & 24 & 2.89256 \\
16 & 16.00 & 32 & 4.43339 \\
\bottomrule
\end{tabular}
\end{table}

For $n=8$, the largest two eigenvalues are $7.9681$ and $7.5752$, accounting
for $97.1\%$ of the trace. This is the finite-window geometry in
Fig.~\ref{fig:geometry-main}; the full-pair calculation uses no small-phase approximation.
\subsection{Two static ordering counterexamples}
\label{sec:static-counterexamples}

Two explicit examples show how phase coupling and interval length change
the ordering of frequency tables.

\paragraph{Cosine-only collision can choose the lower-rank table.}
Take $K=4$, $L=8\pi$, and the two descending frequency tables
\[
\Omega_A=(1,5/8,3/8,1/4),\qquad
\Omega_B=(1,6.01/8,4.01/8,1/4).
\]
They share both endpoints. Define $C_{\cos}$ as the mean squared normalized
cosine--cosine overlap over distinct pairs. For $\Omega_A$, every cosine is
an integer harmonic on $[0,L]$, so $C_{\cos}(A)=0$ exactly. Mixed harmonic
parities nevertheless produce nonzero sine--cosine cross terms. For
$\Omega_B$, the small perturbation of an even-harmonic table gives
\begin{align}
C_{\cos}(A)&=0 < 1.38725\times10^{-5}=C_{\cos}(B),\\
r_2(A)&=6.29478 < 7.99971=r_2(B).
\label{eq:cosine-order-counterexample}
\end{align}
Thus a strict cosine-only preference can have the opposite full-subspace
rank ordering. Both values follow directly from Eq.~\eqref{eq:cross-gram}.

\paragraph{Full-subspace collision can reverse across lengths.}
Now use $K=4$, $L=2\pi$, and the fixed tables
\[
\Omega_A=(1,0.99,0.02,0.01),\qquad
\Omega_B=(1,2/3,1/3,0.01).
\]
Again the endpoints agree. Writing $C_M$ for mean phase-invariant collision
under the uniform measure on $[0,M]$, Eq.~\eqref{eq:cross-gram} gives
\begin{align}
C_L(A)&=0.538996 < 0.631721=C_L(B),\\
C_{2L}(A)&=0.384445 > 0.204995=C_{2L}(B),\\
C_{4L}(A)&=0.345188 > 0.033506=C_{4L}(B).
\label{eq:length-order-counterexample}
\end{align}
The close pairs in $A$ retain high overlap as the interval grows, while the
more evenly spread frequencies in $B$ become easier to distinguish. The
example makes the interval dependence of the geometry explicit.

All frequency entries and measures are specified above. The accompanying
\path{figs/verify_explicit_geometry.py} evaluates the closed-form Gram and
checks it against independent $128$-point Gauss--Legendre quadrature;
the largest discrepancy in the displayed metrics is below $4\times10^{-14}$.
These finite basis examples complement the trained allocation comparisons and
the weight--table crossings.

\FloatBarrier
\subsection{Integer-position kernel equivalence and its boundaries}
\label{sec:discrete-kernel-proof}
\begin{corollary}[Integer-position equivalence]\label{cor:discrete-kernel}
Let $\Omega,\Lambda$ contain $K$ frequencies in $(0,\pi)$ and
$\mathcal R_\Omega(d)=\operatorname{diag}_k R(\omega_kd)$.
Position-independent invertible $A,B$ satisfy
$A^\top\mathcal R_\Omega(d)B=\mathcal R_\Lambda(d)$ for $d=0,1$
if and only if the frequency multisets agree.
\end{corollary}

Corollary~\ref{cor:discrete-kernel} applies to the full bilinear kernel on the
$2K$-dimensional rotary space. Equality for all $q,k$ is matrix equality.
At $d=0$, $A^\top B=I$; at $d=1$ the two rotation matrices are similar.
Their complex eigenvalues are $e^{i\omega_k},e^{-i\omega_k}$ including
multiplicities. On $(0,\pi)$ the positive-angle eigenvalue identifies each
frequency uniquely, proving necessity even with repeated frequencies.
For sufficiency choose a block permutation $P$ with
$\mathcal R_\Lambda(d)=P\mathcal R_\Omega(d)P^\top$, and set
$A=B=P^\top$. This holds for every integer $d$, including negative distances.


The interval $(0,\pi)$ selects unique positive rotation angles: outside it,
$\omega$ and $\omega+2\pi$ are indistinguishable at integer positions.
The statement concerns equality of the full bilinear kernel for every content
pair. These conditions are verified by the bundled discrete-kernel checker.
\subsection{Positional overlap preserves content coordinates}
\label{sec:content-coordinate-retention}
For a block $\mathcal B$ of $r$ rotary pairs, write
$\mathcal R_{\mathcal B}(d)=\operatorname{diag}_{k\in\mathcal B}R(\omega_kd)$.
Each block is orthogonal, so $\mathcal R_{\mathcal B}(d)$ has rank $2r$ at
every distance. This is the rank of an operator on content coordinates;
the positional Gram instead measures shared dependence on distance.

\begin{proposition}[Content matching in a slow block]
\label{prop:slow-content-matching}
Let $H>0$ and $\eta=H\max_{k\in\mathcal B}|\omega_k|$. For $|d|\le H$,
\begin{equation}
\left|q_{\mathcal B}^{\top}\mathcal R_{\mathcal B}(d)k_{\mathcal B}
-q_{\mathcal B}^{\top}k_{\mathcal B}\right|
\le\eta\|q_{\mathcal B}\|_2\|k_{\mathcal B}\|_2.
\label{eq:slow-content-bound}
\end{equation}
For $2r$ content symbols represented by distinct coordinate unit vectors,
let the query and correct key share a vector and all incorrect keys use
other vectors. Their unscaled score margin is at least $1-2\eta$ across
arbitrary key distances in $[-H,H]$.
\end{proposition}

\begin{proof}
For each pair,
$\|R(\omega d)-I\|_2=2|\sin(\omega d/2)|\le|\omega d|$.
The block-diagonal norm is the maximum pair norm. The bilinear norm bound
gives Eq.~\eqref{eq:slow-content-bound}. Unrotated matching scores are one
for the correct symbol and zero for an incorrect symbol. Each score
changes by at most $\eta$, giving the margin bound. In particular,
$\eta<1/2$ preserves the correct ranking.
\end{proof}

As distinct frequencies approach zero, their positional subspaces approach
$\operatorname{span}\{1,d\}$ while the content construction remains
distinct in all $2r$ coordinates. This provides a concrete distinction
between repeated positional dependence and repeated content information,
consistent with prior analyses of frequency use \citep{barbero2025round}.
The construction concerns specified content vectors and the small-phase
regime; the next subsection evaluates the complete finite frequency tables.

\subsection{Task gains with lower full-pair effective rank}
\label{sec:allocation-rank-quality}
We use the Llama mathematical native grid
$\omega_k^N=500000^{-k/64}$, native length $8192$, $s=4$ and band $[18,35]$.
The TailSpline and MrPro formulas determine all frequencies without model
observations. Effective rank uses the full-pair Gram from
Eq.~\eqref{eq:cross-gram}, with uniform separations on $[0,H]$.

\begin{table}[ht]
\centering\small
\caption{\textbf{A successful allocation can have lower positional effective rank.}
Ranks describe the full mathematical frequency tables; task scores reuse
the clean paired evaluations in Table~\ref{tab:clean-length-main}.
The $16$K and $32$K task panels contain $650$ and $2600$ paired prompts.}
\label{tab:allocation-rank-quality}
\begin{tabular}{@{}lrrrr@{}}
\toprule
Window & $r_2$ TailSpline & $r_2$ MrPro & RULER TailSpline (\%) & RULER MrPro (\%)\\
\midrule
$16$K & $8.283391$ & $8.739887$ & $86.0949$ & $82.7051$\\
$32$K & $10.077270$ & $10.214888$ & $68.2660$ & $56.5436$\\
\bottomrule
\end{tabular}
\end{table}

Continuous and integer-position Grams, FP32 rounding and causal-separation
weights retain this rank ordering at $8/16/32$K. The table separates positional
structure from task quality; a lower rank is not used as a selection rule.

\subsection{From distance response to content competition}
\label{sec:content-competition}
Fix the pre-RoPE query and keys for one attention read. With content coefficients
$C_{jk}$ and separations $d_j$, its logits can be written
$\ell_j(\nu)=\operatorname{Re}\sum_k C_{jk}e^{i d_j\nu_k}$, absorbing the
attention normalization into $C_{jk}$. Changing the table gives the finite shift
$\eta_j=\operatorname{Re}\sum_k C_{jk}(e^{i d_j\nu'_k}-e^{i d_j\nu_k})$.
For nonempty target and competing groups $G,B$ partitioning the valid keys,
let $p_G,p_B$ be their separately normalized original softmax weights. The
log-odds $\mathcal O=\log\sum_Ge^{\ell_j}-\log\sum_Be^{\ell_j}$ obey
\begin{equation}
\mathcal O'-\mathcal O
=\log\mathbb E_{p_G}e^{\eta_j}-\log\mathbb E_{p_B}e^{\eta_j}.
\label{eq:content-competition}
\end{equation}
Indeed, factor each original group partition function from its updated sum.
If $\min_G\eta_j>\max_B\eta_j$, the target share strictly increases.
This sufficient condition compares target evidence with its competitors,
rather than using positional rank alone. The identity concerns one fixed-state
read; a full forward pass also changes states in subsequent layers.
It is an explanatory softmax identity, not a task-optimality criterion.
